\documentclass[10pt,a4paper]{amsart}
\usepackage[margin=19mm]{geometry}
\usepackage{amsmath,amssymb,mathtools}
\usepackage[scr=rsfs,cal=euler]{mathalfa}
\usepackage{xfrac}
\usepackage[unicode,hidelinks,bookmarks=false]{hyperref}
\newcommand{\ie}{i.e.\ }
\newcommand{\cf}{cf.\ }
\newcommand{\resp}{respectively\ }
\newcommand{\Subsection}{\subsection}
\providecommand{\nobreakdash}{\nobreak\hbox{-}}
\setlength{\emergencystretch}{2em}
\multlinegap0pt
\usepackage{amssymb}
\newtheorem{proposition}{Proposition}
\def\FF{\mathbb{F}}
\title{Matrices over finite fields of characteristic 2 \\ as sums of diagonalizable and square-zero matrices}
\author{Peter Danchev, Esther Garc\'\i a, Miguel G\'omez Lozano}
\date{Source: arXiv:2604.15286v1 (CC BY 4.0)}
\begin{document}
\maketitle

\noindent\emph{Source and licence.} Matrices over finite fields of characteristic 2 \\ as sums of diagonalizable and square-zero matrices. Version arXiv:2604.15286v1 (CC BY 4.0), distributed by arXiv under the Creative Commons Attribution 4.0 International licence, http://creativecommons.org/licenses/by/4.0/, as recorded in the arXiv licence metadata for this exact version and shown on the abstract page https://arxiv.org/abs/2604.15286v1. Full licence notice: \texttt{licenses/arxiv-2604.15286-CC-BY-4.0.txt}.\par\smallskip

\noindent\textit{Editorial scope.} Counted unit: the article's proposition 4k (source0.tex lines 531--537). The article's complete original proof of it is delivered with it (source0.tex lines 539--759, one \texttt{proof} environment of 221 lines). No mathematics is written, repaired or re-derived: the statement and the proof are the article's own lines, in the article's own notation, and no retained line is substituted or re-flowed.  No further source-local context is needed: every cross-reference of every retained line resolves inside the two delivered spans.  Nothing was omitted from any retained span, so the package carries no omission record.  No retained line contains a citation, so the wrapper carries no bibliography and no bibliography entry.  No retained line contains a figure environment, an image-inclusion command or drawing code, so no image is delivered and the package declares no asset dependency.  Editorial formatting only, in addition: an AMS \texttt{amsart} preamble that restates the article's own declarations and macros used by the retained lines, namely the environments \texttt{proposition} on separate counters, together with the standard packages those lines need.  The retained environments print in this document as Proposition 1 in order of appearance, whereas the article prints the same items with the numbering of its own full document.  The complete native source of the article as distributed in the arXiv e-print for this exact version is bundled unchanged as \texttt{sources/198596/source0.tex}.
\begin{proposition}\label{4k}
Let $\FF$ be a field of characteristic 2 with $\#\FF\ge4$, and let $B$ be a non-derogative matrix over $\FF$ of order $4k$ with $k\ge 2$. Then, for every $a\in \FF$ with $a\ne 0,1$, we have:
  \begin{itemize}
\item[(a)] If $\text{Trace}(B)=c\ne 0$, there exists a square-zero matrix $N\in {M}_{4k}(\FF)$ such that $B=N+D$, where $D$ is diagonalizable with eigenvalues $0, c, ca,c(a+1)$.
\item[(b)] If $\text{Trace}(B)=0$, there exists a square-zero matrix $N\in {M}_{4k}(\FF)$ such that $B=N+D$, where $D$ is diagonalizable with eigenvalues $0,1,a$ and $a+1$.
  \end{itemize}
\end{proposition}

\begin{proof}
{\bf Case (a):} Suppose that $\text{Trace}(B)=c\ne 0$. Without loss of generality, let us replace $B$ by $A=c^{-1}B$, which is again a non-derogative matrix which possesses trace equal to 1; thereby, we can suppose without loss of generality that $A$ has the form
$$
A=\left(\begin{array}{cccccc}
                                                                          0 & 0 & \cdots & 0 & 0 &u_0\\
                                                                          1 & 0 & \cdots & 0 & 0 &u_1\\
                                                                          0 & 1 & \cdots & 0 & 0 &u_2\\
                                                                          \vdots & \vdots & \ddots & \vdots & \vdots &\vdots\\
                                                                          0 & 0 & \cdots & 1 & 0 &u_{4k-2}\\
                                                                          0 & 0 & \cdots & 0 & 1 &1
                                                                        \end{array}\right).$$
We are now going to build a square-zero matrix $N$ such that $A=N+D$, where $D$ can be expressed as the direct sum of $k-2$ diagonalizable blocks of order $4$, two diagonalizable blocks of order $3$ and one diagonalizable block of order $2$. To that goal, let us consider
$$
M_s:=\left(
   \begin{array}{cc}
      0 &u_{4s-4}\\
      0 &u_{4s-3}\\
      0 &u_{4s-1}+u_{4s-2}\\
      0 &0\\
   \end{array}
 \right) \hbox{ for  $s=1,\dots, k-2$,}
$$
\small
 $$M_{k-1}:=\left(
   \begin{array}{cc}
      0 &u_{4k-8}\\
      0 &u_{4k-6}+u_{4k-7}\\
      0 &0\\
   \end{array}
 \right),\  M_{k}:=\left(
   \begin{array}{cc}
      0 &u_{4k-5}\\
      0 &u_{4k-3}+u_{4k-4}\\
      0 &0\\
   \end{array}
 \right),\ M_{k+1}:=\left(
                  \begin{array}{cc}
                    0 & u_{4k-2} \\
                    0 & 0 \\
                  \end{array}
                \right),$$\normalsize
and put
$$N:=\left(
      \begin{array}{cccccccc}
        N_4  & 0 &\cdots & 0& 0 & 0 & 0&  M_1 \\
        Q_{4,4}  &N_4 & \cdots & 0 & 0 & 0& 0  & M_2 \\
        \vdots &\ddots & \ddots & \vdots& \vdots & \vdots& 0  & \vdots \\
        0 &0 & \ddots & N_4 & 0& 0 & 0&  M_{k-3} \\
        0 &0 & \cdots & Q_{4,4} & N_4& 0& 0 &  M_{k-2} \\
        0 &0 & \cdots & 0 & Q_{3,4}& N_3& 0 &  M_{k-1} \\
        0 &0 & \cdots & 0 & 0& Q_{3,3} &N_3&   M_k \\
        0 &0 & \cdots & 0 &0 & 0& Q_{2,3} &  M_{k+1} \\
      \end{array}
    \right)\in M_{4k}(\mathbb{F}).
$$
It can be verified that $N_4 M_s= Q_{4,4} M_s= M_{3,4} Q_s=0$ for $s=1\dots, k-2$, $N_3 M_{k-1}=N_3 M_{k}=Q_{3,3} M_{k-1}= Q_{2,3} M_k=0$, and $M_{k+1}M_{k+1}=0$, so we deduce that $N$ is a square-zero matrix.

Moreover, if we denote
$$T_s:=\left(
   \begin{array}{cc}
       0 &0\\
       0 &0\\
       0 &u_{4s-1}\\
       0 &u_{4s-1}\\
   \end{array}
 \right),\hbox{ for $s=1\dots, k-2$,}
$$
$$T_{k-1}:=\left(
   \begin{array}{cc}
        0 & 0\\
        0 &u_{4k-6}\\
        0 &u_{4k-6}\\
   \end{array}
 \right),\ T_{k}:=\left(
   \begin{array}{cc}
        0 & 0\\
        0 &u_{4k-3}\\
        0 &u_{4k-3}\\
   \end{array}
 \right),
 $$
we have that $A=N+D$, where
$$D:=\left(
         \begin{array}{cccccc}
           D_4 & \cdots & 0 & 0 & 0 & T_1 \\
           \vdots & \ddots & \vdots & \vdots & \vdots & \vdots \\
           0 & \cdots & D_4 & 0 & 0 & T_{k-2} \\
           0 & \cdots & 0 & D_3 & 0 & T_{k-1} \\
           0 & \cdots & 0 & 0 & D_3 & T_k \\
           0 & \cdots & 0 & 0 & 0 & D_2 \\
         \end{array}
       \right).$$
Furthermore, since $D_4$ is diagonalizable, for any $s=1\dots, k-2$, the vectors
\begin{align*}
 v_{4s-3}&=(a^2+a)e_{4s-3}+(a^2+a+1)e_{4s-2}+e_{4s},\\
 v_{4s-2}&=(a^2+a)e_{4s-2}+e_{4s-1}+e_{4s},\\
 v_{4s-1}&=(a+1)e_{4s-2}+a e_{4s-1}+e_{4s},\\
 v_{4s}&=a e_{4s-2}+(a+1) e_{4s-1}+e_{4s}
\end{align*}
form a set of $4k-8$ independent eigenvectors of $D$ associated to the eigenvalues $0,1,a,a+1$. Moreover, since $D_3$ is diagonalizable for $r=0,1$ the vectors
\begin{align*}
 v_{4k-7+3r}&=(a^2+a)e_{4k-7+3r}+e_{4k-6+3r}+e_{4k-5+3r},\\
 v_{4k-6+3r}&=(a+1)e_{4k-7+3r}+a e_{4k-6+3r}+e_{4k-5+3r},\\
 v_{4k-5+3r}&=a e_{4k-7+3r}+(a+1) e_{4k-6+3r}+e_{4k-5+3r}\\
\end{align*}
form a set of $6$ independent eigenvectors of $D$ associated to the eigenvalues $1,a,a+1$, and the following vectors are eigenvectors of $D$ associated to the eigenvalues $0,1$:
\begin{align*}
  v_{4k-1}&=\sum_{i=1}^{k-2}    (u_{4i-1} e_{4i-2}+ u_{4i-1} e_{4i-1})+ \sum_{j=0}^{1} (u_{4k-6+3j} e_{4k-7+3j}+u_{4k-6+3j} e_{4k-6+3j})\\&+    e_{4k-1}+   e_{4k}, \quad  \\
  v_{4k}&=\sum_{i=1}^{k-2}   ((a^2+a) u_{4i-1} e_{4i-2}+ u_{4i-1} e_{4i})\\& +\sum_{j=0}^{1}((a^2+a) u_{4k-6+3j} e_{4k-7+3j} +u_{4k-6+3j} e_{4k-5+3j})+ e_{4k}.
\end{align*}
Consequently, $\{v_1,\dots, v_{4k}\}$ is a basis of eigenvectors of $D$, and hence $D$ is diagonalizable. Observe in addition that, due to the structure of the matrix
$D$, it is readily verified that $v_{4k-1}$ and $v_{4k}$ are eigenvectors associated to $0,1$ within a $9\times 9$ matrix
$$D':=\left(
  \begin{array}{ccc}
    D_4 &0& T_1 \\
    0 & D_3& T_k \\
     0 & 0& D_2 \\
  \end{array}
\right).
$$
We thus have expressed $A$ as the sum of a square-zero matrix and a diagonalizable matrix with eigenvalues $0,1,a,a+1$, whence $B=cA$ can also be expressed as the sum of a square-zero matrix and a diagonalizable matrix with eigenvalues $0,c,ca, c(a+1)$.

\medskip

{\bf Case (b):} Suppose that $\text{Trace}(B)= 0$. Suppose also without loss of generality that $B$ has the form
$$
B=\left(\begin{array}{cccccc}
                                                                          0 & 0 & \cdots & 0 & 0 &u_0\\
                                                                          1 & 0 & \cdots & 0 & 0 &u_1\\
                                                                          0 & 1 & \cdots & 0 & 0 &u_2\\
                                                                          \vdots & \vdots & \ddots & \vdots & \vdots &\vdots\\
                                                                          0 & 0 & \cdots & 1 & 0 &u_{4k-2}\\
                                                                          0 & 0 & \cdots & 0 & 1 &0
                                                                        \end{array}\right).$$
We are now going to build a square-zero matrix $N$ such that $B=N+D$ and $D$ can be decomposed as the direct sum of one diagonalizable block of order $1$, $k-1$ diagonalizable blocks of order $4$ and one diagonalizable block of order $3$: To that end, let us consider
$$M_s:=\left(
   \begin{array}{ccc}
     0 & 0 &u_{4s-3}+(a^2+a) u_{4s}\\
     0 & 0 &u_{4s-2}+(a^2+a+1) u_{4s}\\
     0 & 0 &u_{4s-1}\\
     0 & 0 &0\\
   \end{array}
 \right),\hbox{ for $s=1,\dots, k-1$,}
$$
$$M_{k}:=\left(
   \begin{array}{ccc}
      0& 0 &u_{4k-3}+a^2+a\\
      0 & 0 &u_{4k-2}+a^2+a+1\\
      0 & 0 &0\\
   \end{array}
 \right)$$
and put
$$N:=\left(
      \begin{array}{cccccccc}
        0 & 0 & 0 & \cdots & 0 &0 &  0 \\
        Q_{4,1} & N_4 & 0 & \cdots &0 & 0  & M_1 \\
        0 & Q_{4,4} & N_4 & \cdots &0 & 0  & M_2 \\
        \vdots & \vdots & \ddots & \ddots  & \vdots & \vdots \\
        0 & 0 & 0 & \ddots &N_4 & 0 &  M_{k-2} \\
        0 & 0 & 0 & \cdots &Q_{4,4} & N_4 &  M_{k-1} \\
        0 & 0 & 0 & \cdots &0 & Q_{3,4} &  M_k \\
      \end{array}
    \right)\in M_{4k}(\mathbb{F}).
$$
It can be verified that $N_4 M_s=Q_{4,4} M_s$ for $s=1,\dots, k-1$, $Q_{3,4}M_{k-1}=0$, and $M_{k} M_{k}=0$, we conclude that $N$ is a square-zero matrix. Moreover, if we denote
$$
T_s:=\left(
   \begin{array}{ccc}
      0 & 0 &(a^2+a)u_{4s}\\
      0 & 0 &(a^2+a+1)u_{4s}\\
      0 & 0 &0\\
      0 &  &u_{4s}\\
   \end{array}
 \right),\hbox{$s=1\dots, k-1$, and}
$$
 $$T_0:=\left(
                 \begin{array}{ccc}
                   0 & 0 & u_0 \\
                 \end{array}
               \right),\quad T_{k}:=\left(
   \begin{array}{ccc}
       0 & 0 &a^2+a\\
       0 & 0 &a^2+a+1\\
       0 & 0 &0\\
   \end{array}
 \right),
 $$
we obtain that $B=N+D$, where
$$D:=\left(
         \begin{array}{ccccc}
           0 & 0 & \cdots  & 0 & T_0 \\
           0 & D_4 & \cdots  & 0 & T_1 \\
           \vdots & \vdots & \ddots  & \vdots & \vdots \\
           0 & 0 & \cdots  & D_4 & T_{k-1} \\
           0 & 0 & \cdots  & 0 & D_3 \\
         \end{array}
       \right).$$
Additionally, notice that $v_1=e_1$ is a eigenvector of $D$ associated to the eigenvalue $0$. Moreover, since $D_4$ is diagonalizable, for any $s=1,\dots, k-1$ the vectors
\begin{align*}
 v_{4s-2}&=(a^2+a)e_{4s-2}+(a^2+a+1)e_{4s-1}+e_{4s+1},\\
 v_{4s-1}&=(a^2+a)e_{4s-1}+e_{4s}+e_{4s+1},\\
 v_{4s}&=(a+1)e_{4s-1}+a e_{4s}+e_{4s+1},\\
 v_{4s+1}&=a e_{4s-1}+(a+1) e_{4s}+e_{4s+1}
\end{align*}
form a set of $4k-4$ independent eigenvectors of $D$ associated to $0,1,a,a+1$, and  the following vectors are eigenvectors of $D$ associated to the eigenvalues $1,a,a+1$:
\begin{align*}
  v_{4k-2}&=u_0 e_1+\sum_{i=1}^{k-1}  ( (a^2+a) u_{4i} e_{4i-2}+ (a^2+a+1) u_{4i} e_{4i-1}+ u_{4i} e_{4i+1})\\&+ (a^2+a) e_{4k-2}+  e_{4k-1}+ e_{4k},\\
  v_{4k-1}&=u_0 e_1+\sum_{i=1}^{k-1}  ( (a^2+a) u_{4i} e_{4i-2}+ a^2 u_{4i} e_{4i-1}+a u_{4i} e_{4i})\\&+ (a^2+a) e_{4k-2}+a^2  e_{4k-1}+ a e_{4k},\\
  v_{4k}&=u_0 e_1+\sum_{i=1}^{k-1}  ( (a^2+a) u_{4i} e_{4i-2}+ (a^2+1) u_{4i} e_{4i-1}+(a+1) u_{4i} e_{4i})\\&+ (a^2+a) e_{4k-2}+(a^2+1)  e_{4k-1}+ (a+1) e_{4k}.\\
\end{align*}
Therefore, $\{v_1,\dots, v_{4k}\}$ is a basis of eigenvectors of $D$, and hence $D$ is diagonalizable. Observe in addition that, due to the structure of the matrix
$D$, it is easily verified that $v_{4k-2}, v_{4k-1}$ and $v_{4k}$ are eigenvectors associated to the eigenvalues $1,a,a+1$ within a $7\times 7$ matrix
$$D':=\left(
  \begin{array}{cc}
    D_4 & T_1 \\
    0 & D_3\\
  \end{array}
\right).
$$
We thus have expressed $B$ as the sum of a square-zero matrix and a diagonalizable matrix with eigenvalues $0,1,a,a+1$, as promised.
\end{proof}

\end{document}
